\section{Conclusions and limits}\label{sec:conclusions}

\paragraph{Conclusions.}
\begin{itemize}[nosep]
\item \textbf{The real noise is the main result.} Whitened O1 H1 data are Gaussian in the bulk, but a glitch floor of a few $10^{-3}$ per second and strong epoch-to-epoch variations decide the price of a false-alarm probability: lowering it from $10^{-2}$ to $10^{-3}$ moves the 50\,\% point of the bank from $\snr\approx5.5$ to $\approx11$ (in Gaussian noise the threshold would move only from $\approx27$ to $\approx31$, i.e.\ by $\approx0.4$ in SNR), and a threshold calibrated on one set of epochs is off by a factor $\approx4.5$ on others.
\item \textbf{Frequentist.} The profile-likelihood bank reaches 50\,\% efficiency at $\snr\approx5$ (oracle with known parameters: $3.0$); the gap is the look-elsewhere penalty, quantified by $N_{\rm eff}\approx6800$ independent trials. The MLE saturates the Fisher bound for $\snr\gtrsim10$; the amplitude estimate is biased and the profile interval under-covers near the boundary, where Neyman/Feldman--Cousins intervals are exact. The dominant analysis systematic is the PSD estimator (mean against median Welch: $+0.025$ on $\varepsilon(\snr=8)$).
\item \textbf{Bayesian.} The exact grid posterior is validated against a closed form, MCMC and simulation-based calibration, in Gaussian and in real noise; the Bayes factor ranks windows like $\lmax$ but is prior-dependent (Occam factor $-2\ln\sigma_\rho$) and exceeds 100 on $0.2$--$0.9\,\%$ of pure-noise windows because of glitches. Credible intervals match confidence intervals for $\snr\gtrsim3$ and differ, depending on the prior, below.
\item \textbf{Machine learning.} A small CNN ensemble matches the bank on sine-Gaussians and beats it on the pulses and chirps it was trained on by about one SNR unit, but a network trained on one morphology fails on others; an autoencoder trained on noise adds nothing to an energy statistic. Calibration is good ($\mathrm{ECE}\lesssim3\,\%$), mostly limited by the epoch shift, and ensemble uncertainty identifies the windows where the model should not be trusted.
\item \textbf{Sequence models and tuning.} A Transformer and a GRU detect worse than the CNN, with the same data and with every lever tried; the training loss of the Transformer stays high ($\approx0.3$ against $0.12$ for the CNN): the patch embedding cannot build the frequency-selective filters a burst needs. A convolutional stem in front of the Transformer recovers most of the gap but does not overtake the CNN. The largest effect of the whole study is on the CNN itself, which was data-limited: fresh injections at every epoch make a single network better than the five-member ensemble ($\overline{\snr_{50}}=4.6$ against $4.85$ at FAP $10^{-2}$).
\item \textbf{Which morphology.} Bayesian model comparison is the best classifier ($0.97$ three-way accuracy at $\snr=8$), the CNN is tied with it at high SNR; no method flags a burst that belongs to none of the families.
\end{itemize}

\paragraph{Limits.}
\begin{itemize}[nosep]
\item \textbf{Single detector.} No coincidence and no time slides: the false-alarm numbers are per one-second window, not rates, and the glitch floor would be largely removed by a second detector.
\item \textbf{Injections at a given SNR.} Not converted into distances or energies; the three morphologies are phenomenological and share the same $t_0$ prior.
\item \textbf{Background statistics.} $1.6\times10^4$ training windows limit the lowest usable FAP to $\sim10^{-3}$; the quoted $10^{-3}$ efficiencies carry a large uncertainty from the tail of the background.
\item \textbf{Eight chunks.} The epoch dependence is measured with two test chunks; a systematic study of its size would need many more.
\item \textbf{Grid resolution.} The Bayesian grid under-resolves $f_0$ and $t_0$ for $\snr\gtrsim15$ (the SBC coverage of those two parameters, $0.50$--$0.60$, reflects this, not the model).
\item \textbf{ML numbers} are efficiencies \emph{for the training injection prior}; the supervised detectors are only as good as the injections they have seen, and parameter estimation with ML was not attempted.
\item \textbf{Tuning study.} Single seed for the Transformer, GRU and hybrid variants, a budget of 12 epochs, no sweep of depth or width; the multiclass networks were trained on static data and probably under-estimate a CNN trained with fresh injections.
\item \textbf{Comparison thresholds.} The bank is compared with the ML detectors on the validation background; the choice of background changes the numbers by up to $0.6$ in $\snr_{50}$ (Section~\ref{sec:ml}).
\end{itemize}

\paragraph{Natural extensions.} A second detector (L1) for coincidence and a real time-slide background; a rolling background close in time to the data; fine grids or nested sampling for the Bayesian part; training the networks with a glitch class (Gravity-Spy style) and with astrophysical waveforms; parameter regression with the CNN compared with the posterior of Section~\ref{sec:bayes}; a second seed for the Transformer variants and fresh injections for the multiclass networks; an open-set or goodness-of-fit head for bursts outside the three families; longer windows, where attention could matter.
